\chapter{Radiofrequency Neurotomy}

\section*{What Is This Test or Procedure?}
Radiofrequency neurotomy (rhizotomy) is a minimally invasive procedure that uses heat generated by a radiofrequency current to interrupt the nerve supply to a painful joint or structure. It is performed under fluoroscopic or CT guidance to deliver targeted lesions to specific medial branch nerves or other sensory nerves.

\section*{Alternative Names}
Radiofrequency ablation of nerves, radiofrequency rhizotomy, medial branch radiofrequency neurotomy, RF neurotomy, facet denervation, pulsed radiofrequency therapy

\section*{Principle}
A needle electrode is advanced under image guidance to a target nerve, and a radiofrequency current heats the surrounding tissue to denature nerve proteins and produce a thermal lesion. This interrupts transmission of nociceptive signals from the affected joint or structure to the spinal cord. Because the lesion is confined to the target nerve, adjacent motor function is typically preserved. The procedure is diagnostic in concept when combined with temporary blocks that predict pain relief from subsequent denervation.

\section*{History}
Percutaneous radiofrequency lesioning of spinal nerves was introduced in the 1970s as a development of earlier surgical rhizotomy techniques. It became widely used in pain medicine for cervical and lumbar facet arthropathy after studies validated its efficacy and low complication rate.

\section*{Why This Test or Procedure Is Ordered}
Radiofrequency neurotomy is ordered for chronic axial neck or back pain attributed to facet joint arthropathy that has not responded to conservative management and that is confirmed by positive responses to diagnostic medial branch blocks. It may also be used for sacroiliac joint pain, occipital neuralgia, and selected cases of discogenic or postsurgical pain. The goal is durable, long-lasting pain relief and functional improvement when injectable therapies provide only temporary benefit.

\section*{Is This a Primary or Secondary Test or Procedure}
It is a secondary, confirmatory procedure because it follows diagnostic blocks and is reserved for patients with a clear pain generator. It is not a primary test but a definitive therapeutic intervention for selected chronic pain conditions.

\section*{How This Test or Procedure Is Performed}
The patient lies prone, and the target medial branch nerves are identified with fluoroscopy or CT, then anesthetized locally. A radiofrequency cannula with electrode is advanced to each target site, and sensory and motor stimulation is performed to confirm correct placement and avoid motor nerves. After test stimulation, thermal lesions are created at each target nerve for approximately 60--90 seconds each, and the needles are removed.

\section*{What Preparation Is Needed}
Patients stop anticoagulant and antiplatelet medications before the procedure, and a coagulation panel may be checked. No special fasting is usually required unless sedation is planned. Imaging and a review of prior diagnostic blocks are used to plan the level and number of targets.

\section*{Contraindications and Precautions}
Uncorrectable coagulopathy, active infection at the site, and pacemakers or implanted devices that may be affected by radiofrequency current are contraindications. The procedure should not be performed without a positive diagnostic block, and patients with widespread or non-mechanical pain may not benefit. Care must be taken near motor nerves, the cervical region, and the vertebral artery.

\section*{Risks}
Risks are generally low and include transient procedural pain, bruising, bleeding, infection, and neuritis with temporary burning or dysesthesia. Rare complications include incomplete pain relief, motor weakness, and injury to adjacent neurovascular structures, especially at cervical levels.

\section*{What Happens After the Test or Procedure}
Patients are observed briefly and usually return home the same day with instructions to limit strenuous activity for a few days. Postprocedural soreness is expected and managed with cold packs and analgesics. Pain relief, when it occurs, typically develops over 1--2 weeks as the nerve lesion matures.

\section*{Understanding Your Results}
A successful neurotomy produces significant pain reduction that may last 6--12 months or longer as the treated nerves regenerate slowly. If pain returns, the procedure can be repeated, and recurrence at the same level often responds equally well. Lack of relief after a technically successful neurotomy suggests the pain source is not the targeted nerve.

\section*{Normal Results}
There is no normal numeric result because neurotomy is a therapeutic intervention; the expected outcome is marked pain reduction and improved function. Success is typically defined as at least 50\%--80\% pain relief for a sustained period.

\section*{Follow-up Tests and Procedures}
Patients are followed clinically to gauge the duration of relief, and repeat neurotomy may be offered when pain recurs. If relief is inadequate, further imaging or additional diagnostic blocks may be considered to reevaluate the pain generator.

\section*{References}
\begin{enumerate}
  \item MedlinePlus. Radiofrequency Neurotomy. U.S. National Library of Medicine; 2025.
  \item Current Medical Diagnosis and Treatment. McGraw-Hill; 2025.
\end{enumerate}

\clearpage
